\documentclass[12pt]{report}

% Packages
\usepackage[
    left=1.5in,
    right=1.25in,
    top=1.25in,
    bottom=1.25in
]{geometry}

\usepackage{setspace}          % line spacing
\usepackage{graphicx}          % figures
\usepackage{amsmath}           % math
\usepackage{hyperref}          % clickable refs
\usepackage{float}             % figure placement [H]
\usepackage{booktabs}          % clean tables
\usepackage{caption}           % caption formatting
\usepackage{times}             % Times New Roman font
\usepackage{fancyhdr}          % header/footer control
\usepackage{emptypage}         % blank pages have no header/footer
\usepackage{listings}          % code listings
\usepackage{xcolor}            % colors for listings

% Code listing style
\lstset{
    basicstyle=\small\ttfamily,
    breaklines=true,
    frame=single,
    numbers=left,
    numberstyle=\tiny,
    keywordstyle=\bfseries,
    commentstyle=\itshape\color{gray},
}

% Page numbering style
\pagestyle{fancy}
\fancyhf{}
\fancyfoot[C]{\thepage}
\renewcommand{\headrulewidth}{0pt}

% Double spacing (required by UCSC)
\doublespacing

\begin{document}

% ============================================================
% FRONT MATTER
% ============================================================

% Title page (counted as page i, no number displayed)
\begin{titlepage}
    \centering

    {\large UNIVERSITY OF CALIFORNIA}\\[0.1in]
    {\large SANTA CRUZ}

    \vspace{0.5in}

    {\large \textbf{JAZZYFS: ADAPTIVE PREFETCHING WITH LIGHTWEIGHT\\
    FEEDBACK-DIRECTED SEQUENTIAL PREDICTION\\
    AND POST-WORKLOAD SONIFICATION}}

    \vspace{0.4in}

    {\normalsize A thesis submitted in partial satisfaction of the\\
    requirements for the degree of}

    \vspace{0.3in}

    {\large \textbf{MASTER OF SCIENCE}}

    \vspace{0.2in}

    {\normalsize in}

    \vspace{0.2in}

    {\large \textbf{COMPUTER SCIENCE AND ENGINEERING}}

    \vspace{0.4in}

    {\normalsize by}

    \vspace{0.2in}

    {\large \textbf{Michelle Gurovith}}

    \vspace{0.4in}

    {\normalsize June 2026}

    \vspace{0.6in}

    \begin{flushleft}
    The Thesis of Michelle Gurovith is approved:\\[0.3in]
    \rule{3in}{0.4pt}\\
    Professor Scott Brandt, Chair\\[0.3in]
    \rule{3in}{0.4pt}\\
    Professor Andi Quinn\\[0.3in]
    \rule{3in}{0.4pt}\\
    Professor Lise Getoor\\[0.3in]
    \rule{3in}{0.4pt}\\
    Professor Heiner Litz\\[0.5in]
    \rule{3in}{0.4pt}\\
    Peter Biehl\\
    Vice Provost and Dean of Graduate Studies
    \end{flushleft}

\end{titlepage}

% Copyright page (counted as page ii, no number displayed)
\thispagestyle{empty}
\vspace*{\fill}
\begin{center}
    Copyright \copyright\ 2026 by Michelle Gurovith\\
    All rights reserved.
\end{center}
\vspace*{\fill}
\newpage

% Front matter uses Roman numerals starting at iii
\pagenumbering{roman}
\setcounter{page}{3}

% Table of Contents
\tableofcontents
\newpage

% List of Figures
\listoffigures
\newpage

% List of Tables
\listoftables
\newpage

% Abstract
\chapter*{Abstract}
\addcontentsline{toc}{chapter}{Abstract}

Modern filesystems use heuristic-based prefetching, which uses the same
read-ahead method no matter what the workload is. This is fine for
sequential access, but it uses up cache space and bandwidth in random or
mixed workloads.

JazzyFS is a FUSE-based filesystem that employs adaptive prefetching with
phase detection and confidence scoring. It looks at access patterns while
the program is running and only prefetches when it is sure that the workload
is sequential. A sonification feature turns system behavior into sound,
offering users a post-workload audio summary of system behavior.

We tested JazzyFS on seven workloads on both Linux with ext4 and macOS with
APFS, including a prefetch depth sweep over depths 1, 2, 4, and 8. In every experimental run, the phase detector identified the expected
access pattern type, a run-level consistency check across all seven
workloads. Deeper prefetch depths did
not improve performance and added more overhead. The FUSE layer itself, not
the adaptive prefetching technique, is the main source of overhead.

\newpage

% Acknowledgments
\chapter*{Acknowledgments}
\addcontentsline{toc}{chapter}{Acknowledgments}

I want to thank my advisor and committee for all the help and advice they
gave me during this project. I appreciate the constructive criticism and
support that helped me choose the direction of my research and write this
paper.

\newpage

% ============================================================
% MAIN MATTER
% ============================================================
\pagenumbering{arabic}
\setcounter{page}{1}

% ============================================================
% Chapter 1: Introduction
% ============================================================
\chapter{Introduction}
\label{ch:introduction}

There is always a latency gap between storage and the processor in modern
computing systems. The CPU has to wait for storage to respond when it wants
data that is not currently in memory, which stops execution. Prefetching
fixes this by putting data into memory before it is needed. But modern
filesystems use heuristic-based prefetching, which uses the same fixed
read-ahead rule no matter what the workload is. This works for workloads
that run in order, but it wastes bandwidth and cache space for workloads
that run in random or mixed order. For instance, a web server may visit
several unrelated files in no predictable sequence, which is common in
modern systems. Heuristic prefetching keeps requesting data even when the
pattern is irregular, wasting I/O resources and filling the cache with data
that is never used, because it cannot observe or adapt to how the workload
is behaving.

\section{Motivation}

The main goal of this thesis is to find out how well an adaptive prefetcher
can discover workload I/O patterns, put them in categories, and make
decisions about turning prefetching on or off. For example, if a filesystem
can identify that the requests are contiguous, it can classify the workload
as sequential and enable prefetching. If, on the other hand, the reads come
in at non-contiguous offsets, the system may not be as sure that this is a
sequential workload.

JazzyFS does not use a fixed rule like a heuristic prefetcher. Instead, it
looks at the workload and figures out which pattern it most likely belongs
to. If the workload has a sequential pattern, it acts like a normal
prefetcher to improve I/O efficiency. If the filesystem sees that the
workload pattern is random, it does not prefetch in order to save resources.

Another motivation is observability. Engineers typically need to use
debugging tools to figure out what prefetching is doing. It can take a lot
of effort to understand what the prefetcher was doing using traditional
tools like logs. JazzyFS solves this by playing back a summary of system behavior as audio
after each workload completes, giving users another way to review what
happened beyond just looking at data logs.

\section{JazzyFS}

To deal with these problems, this thesis presents JazzyFS, a read-only
user-space filesystem that is built on top of Filesystem in Userspace
(FUSE). JazzyFS mounts an existing directory and acts as a filter for every
file read operation. It looks at the recent access history for each read,
uses that information to figure out the access pattern, and then decides
whether or not to prefetch. There are three ways that JazzyFS can prefetch:
a Null Prefetcher, a One Block Lookahead prefetcher, and a
Feedback-Directed Sequential Prefetcher, which only prefetches when the
sliding-window confidence score shows that the access pattern is sequential.
There is also a sonification feature that turns system behavior into music,
which is another way to observe how the system works.

\section{Experimental Evaluation}

We test JazzyFS on two platforms, Linux with ext4 and macOS with APFS,
using seven different workloads. There are twenty repetitions for each
condition, which means each platform has 420 experimental runs across the
three prefetching modes. A native timing baseline collected without FUSE
usage gives a benchmark with no overhead. A separate prefetch depth sweep
across $\{1, 2, 4, 8\}$ blocks adds more runs, bringing the total to
1{,}120 runs per platform across all conditions.

The workloads cover a variety of access patterns:

\begin{itemize}
    \item \textbf{Sequential}: reading a 100~MB file in a straight line.
    \item \textbf{Random}: deterministic seeks at offsets that are not next
    to each other.
    \item \textbf{Phase Change}: sequential access transitioning to random
    and back again.
    \item \textbf{Concurrent}: twenty random reads at the same time.
    \item \textbf{Tar Extraction}: taking 100 small files out of an archive.
    \item \textbf{Python Import}: reading 200 small Python module files.
    \item \textbf{Cache Lookup}: 200 page-aligned random seeks with 10~ms
    delays between reads.
\end{itemize}

Metrics collected include execution time, overhead versus native, prefetch
rate, prediction confidence, phase detection accuracy, and a prefetch depth
sweep across $\{1, 2, 4, 8\}$ blocks.

\section{Research Questions}

This thesis is organized around five research questions that help explore
the behavior of lightweight adaptive prefetching in a real filesystem.

\begin{enumerate}

    \item \textbf{Broader Systems Perspective.} What do the observed results
    suggest about the broader applicability of lightweight adaptive techniques
    to systems resource management?

    \item \textbf{Adaptation Dynamics.} How quickly can a feedback-directed
    prefetcher adapt when access patterns change abruptly?

    \item \textbf{Cost vs.\ Benefit.} Are there scenarios where simple
    heuristics are consistently better than adaptive approaches, and vice
    versa?

    \item \textbf{Confidence and Performance.} How does prediction confidence
    correlate with actual performance outcomes under different workloads?

    \item \textbf{Observability Through Sound.} Does mapping access phase,
    confidence, and prefetching mode to musical parameters produce an auditory
    signal that is meaningful and distinguishable,
    enabling operators to monitor system behavior without consulting logs or
    dashboards?

\end{enumerate}

\section{Key Findings}

\begin{itemize}
    \item In every experimental run on both platforms, the phase detector
    produced the expected access pattern type. This is a run-level
    consistency check: it confirms the detector tracks dominant patterns
    correctly, not that every individual read is correctly classified.
    \item Confidence correctly tracks workload structure: $\approx 0.999$ for
    sequential workloads, 0.000 for purely irregular workloads
    (random, concurrent, cache lookup), and $\approx 0.500$ for mixed-pattern
    workloads such as Python module import.
    \item The adaptive algorithm matches OBL on the sequential workload
    (prefetch rate 1.00) and tar extraction (prefetch rate $\geq 0.97$ across
    platforms), and matches the null prefetcher on purely irregular workloads
    (prefetch rate 0.00).
    \item FUSE overhead dominates algorithm cost: even the null prefetcher
    adds 180\% overhead on macOS and 388\% on Linux for sequential workloads.
    \item Prefetch depth greater than one provides no throughput benefit and
    consistently increases overhead on sequential workloads.
\end{itemize}

\section{Contributions}

\begin{enumerate}
    \item Design and implementation of JazzyFS, a FUSE-based adaptive
    filesystem with cross-platform support (Linux ext4 and macOS APFS) and
    full experimental instrumentation.
    \item A lightweight feedback-directed sequential prefetching mechanism
    requiring no training data, implemented as a sliding-window confidence
    scorer.
    \item A reproducible evaluation framework with seven workloads, three
    primary modes, a prefetch depth sweep across four depths, and a native
    baseline, yielding 2{,}240 total experimental runs across both platforms.
    \item A novel sonification system mapping access phase, confidence, and
    mode to a post-workload musical summary.
    \item Quantitative analysis of trade-offs between heuristic and adaptive
    prefetching across diverse workload types on two real filesystems.
\end{enumerate}


% ============================================================
% Chapter 2: Background
% ============================================================
\chapter{Background}
\label{ch:background}

This chapter provides the technical background necessary to understand
JazzyFS.
Section~\ref{sec:bg-fs} reviews filesystem architecture.
Section~\ref{sec:bg-fuse} introduces FUSE.
Section~\ref{sec:bg-prefetch} surveys classical prefetching techniques.
Section~\ref{sec:bg-feedback} reviews feedback-directed prefetching.
Section~\ref{sec:bg-obs} covers observability and sonification.

\section{Filesystem Architecture}
\label{sec:bg-fs}

A filesystem is the software layer that organizes data on storage and
presents it to applications through a common interface. From an
application's perspective, a file is just a named sequence of bytes. The
filesystem translates operations like open, read, write, seek, and close
into physical reads and writes on the storage device.

An important part of the storage stack is the \emph{page cache}. The kernel
keeps recently accessed file data in memory so that repeated reads can be
served from memory instead of storage. When the requested data is already in
cache, no I/O is needed. When it is not, the kernel issues a physical read.
This reduces latency significantly for files that are accessed often or
read sequentially.

\subsection{Access Patterns and I/O Performance}

Access patterns in filesystem workloads span a diverse spectrum from
sequential to random. In a purely sequential workload, each read request
begins exactly where the previous request ended. On the other hand, random
access patterns are offsets that are not contiguous. While these are two
widely seen workloads, pure sequential and random workloads are rare; real
workloads are often a mix of the two. As a result, a properly functioning
filesystem must be able to handle a wide range of workloads without
assigning mixed workloads to a specific pure pattern.

\subsection{ext4: Linux's Default Filesystem}

ext4 is the default filesystem on most Linux systems and is used for
JazzyFS's Linux evaluation. It is a journaled filesystem, which means it
logs changes before writing them, allowing recovery after a crash. The
Linux kernel includes a built-in read-ahead mechanism that loads data ahead
of sequential reads. When a workload reads sequentially, the kernel detects
this and issues extra reads in advance, hiding storage latency.

JazzyFS uses the \texttt{direct\_io=True} FUSE option, which bypasses the
page cache. This means every read goes through the FUSE callback and cannot
be served from memory. Because of this, the native baseline is always faster
than any FUSE mode: native reads get kernel read-ahead and caching, while
JazzyFS reads pay the full FUSE round-trip cost every time. This is important when reading the overhead numbers in the evaluation.

\subsection{APFS: Apple's Flash-Optimized Filesystem}

APFS (Apple File System) is the default filesystem on macOS and is used for
JazzyFS's macOS evaluation. It was designed for solid-state storage and
replaced the older HFS+ filesystem. Unlike spinning disks, solid-state drives
have low random-access latency, so APFS does not optimize specifically for
sequential reads. APFS also uses a copy-on-write design, meaning changes
create new data blocks rather than overwriting existing ones, which enables
atomic snapshots and fast file cloning.

On macOS, FUSE support is provided by macFUSE, a third-party kernel
extension. Because macFUSE and the Linux FUSE module have different
implementations, the per-operation overhead differs between the two
platforms. The evaluation reports results for both platforms separately.

\section{FUSE: Filesystem in Userspace}
\label{sec:bg-fuse}

FUSE (Filesystem in Userspace) is a framework for implementing a filesystem
as a regular user-space program instead of kernel code. This makes
development and testing easier: a bug crashes the process, not the kernel.
Filesystem logic can be written in any language with FUSE bindings and
deployed without kernel patches or reboots.

When an application reads a file on a FUSE mount, the kernel sends the
request to the user-space filesystem process and waits for the result. This
round-trip happens on every read and adds latency. On Linux, the cost is
typically 10--100 microseconds per operation, which adds up quickly for
workloads with many small reads. macFUSE on macOS adds a similar cost.

JazzyFS uses FUSE because the goal is to study prefetching behavior, not to
build a production filesystem. FUSE provides a clear intercept point on
every read, which is exactly what is needed to observe access patterns and
test different prefetching strategies. The FUSE overhead is measured and
reported in the evaluation so it can be separated from the cost of the
prefetching logic itself.

\subsection{FUSE on Linux and macOS}

FUSE works differently on Linux and macOS. On Linux, FUSE is built into
the kernel as a loadable module (\texttt{fuse.ko}), and the \texttt{fusepy}
library communicates with it through \texttt{/dev/fuse}. On macOS, the
kernel does not include FUSE natively, so macFUSE (formerly OSXFUSE)
provides it as a third-party kernel extension. Because the two
implementations use different kernel paths, their per-operation costs
differ. In the JazzyFS evaluation, the null prefetcher adds approximately
388\% overhead on Linux and approximately 180\% on macOS for sequential
workloads, which reflects this difference.

\section{Prefetching Techniques}
\label{sec:bg-prefetch}

Prefetching loads data into memory before it is requested, reducing the
time the CPU spends waiting on storage. The idea is that if past access
patterns are a good predictor of future ones, loading ahead can hide
latency. When the prediction is correct, performance improves. When it is
wrong, bandwidth is wasted and the cache fills with data that is never
used. JazzyFS uses three strategies that range from no prefetching to
always prefetching to prefetching only when the pattern looks sequential,
which makes it possible to isolate the cost and benefit of each approach.

\subsection{Sequential Prefetching}

The simplest prefetching strategy is sequential read-ahead: after reading
block $b$, fetch block $b+1$ into the page cache. This is built into the
Linux kernel's \texttt{readahead} subsystem. It works well for sequential
workloads, but it has no awareness of access patterns. It always fetches
ahead regardless of whether the workload is actually sequential, which
wastes bandwidth and fills the cache with unused data on irregular
workloads.

\subsection{Fixed-Policy Read-Ahead}

JazzyFS's OBL baseline implements the One Block Lookahead strategy. After
every read at offset $o$ for $n$ bytes, it issues a speculative read at
offset $o + n$ for 4~KB. The policy is unconditional: it fires on every
read regardless of the access pattern. On sequential workloads, this works
well because the prefetched block is almost always the next one needed. On
non-sequential workloads, every prefetch misses, consuming bandwidth and
filling the cache with data that is never used. In the JazzyFS evaluation,
OBL represents the always-prefetch extreme and provides a direct comparison
point for the adaptive strategy.

\section{Feedback-Directed Prefetching}
\label{sec:bg-feedback}

Feedback-directed prefetching improves on static strategies by using the
results of past prefetches to control future ones. The idea is that a
prefetcher can track whether the blocks it fetched were actually used. If
prefetches are consistently hitting, it can fetch more aggressively. If they
are consistently missing, it can slow down or stop. This allows the
prefetcher to adapt to the workload without any manual configuration or
training data.

Traditional feedback-directed prefetchers work at the hardware level and
use performance counters in the CPU or memory controller to measure hit
rates. JazzyFS takes a simpler approach: rather than tracking whether
prefetched blocks were actually used, it uses the contiguity of incoming
read offsets as a proxy signal. If recent reads are contiguous, confidence
is high and prefetching is enabled. If they are scattered, confidence is
low and prefetching is skipped. This is a pattern-observation signal, not
a hit/miss feedback loop. The full mechanism is described in
Chapter~\ref{ch:design}.

\section{Observability in Systems}
\label{sec:bg-obs}

Observability is the ability to understand what a system is doing from its
external outputs. In storage systems, this usually means looking at access
patterns, latency, and cache behavior. Traditional tools like
\texttt{blktrace}, \texttt{perf}, and DTrace provide detailed data, but
they require post-processing. An engineer must collect a trace, then
analyze it separately before drawing any conclusions about what the system
was doing. There is no immediate feedback while a workload is running.

\subsection{Sonification}

Sonification is the use of non-speech audio to convey data or
information. Unlike visualization, which requires
attention to a screen, sonification can operate in the periphery of awareness,
allowing a practitioner to monitor system state while focusing on other work.

Sonification has been applied in medical monitoring, astronomy, and network
security. Its application to filesystem
behavior is less explored. A well-established sonification design principle
is to map independent state dimensions to orthogonal perceptual parameters
so that each audio channel carries a distinct, non-redundant signal.
JazzyFS's sonification design, described in Chapter~\ref{ch:design},
follows this principle.


% ============================================================
% Chapter 3: Design
% ============================================================
\chapter{Design}
\label{ch:design}

This chapter describes the design and implementation of JazzyFS.
Section~\ref{sec:design-arch} presents the overall architecture.
Section~\ref{sec:design-window} describes the sliding-window access tracker.
Section~\ref{sec:design-phase} explains the phase detection algorithm.
Section~\ref{sec:design-confidence} describes the confidence scorer.
Section~\ref{sec:design-strategies} specifies the three prefetching strategies.
Section~\ref{sec:design-sonification} presents the sonification engine.
Section~\ref{sec:design-read} describes the read callback logic.
Section~\ref{sec:design-impl} covers implementation details.

\section{Architecture Overview}
\label{sec:design-arch}

JazzyFS is a passthrough FUSE filesystem. It mounts an existing directory
and forwards all filesystem operations to the underlying filesystem. It is
read-only; write operations are rejected. All prefetching and observation
logic lives on the read path only. The underlying filesystem handles actual
storage, and JazzyFS only intercepts reads to observe access patterns and
decide whether to prefetch.

The main class is \texttt{PassthroughRO}, which implements five FUSE
callbacks: \texttt{getattr}, \texttt{readdir}, \texttt{open}, \texttt{read},
and \texttt{release}. All prefetching and phase detection logic runs inside \texttt{read()}; sonification playback runs in a separate background thread.
The other callbacks pass requests through with no added logic.

\section{Sliding-Window Access Tracker}
\label{sec:design-window}

JazzyFS keeps a sliding window of the last 1000 read events. Each entry
stores the file path, byte offset, and read size. The 1000-entry cap keeps
memory usage bounded no matter how long a workload runs. Phase detection
and confidence scoring only look at the most recent \texttt{PHASE\_WINDOW=5}
entries, so both stay fast regardless of total run length.

The window clears automatically if more than 2 seconds pass between reads.
This separates back-to-back workload runs without needing an explicit reset,
and also handles the case where one workload ends and a new one begins
without remounting.

\section{Phase Detection}
\label{sec:design-phase}

The phase detector looks at a window of up to the last five reads and classifies
the access pattern as \texttt{sequential}, \texttt{irregular}, or \texttt{unknown}.
A transition between two consecutive reads is sequential if the current
read starts exactly where the previous one ended:
\[
\text{prev.offset} + \text{prev.size} = \text{curr.offset}
\]
If all observed transitions in the current window are sequential, the phase is
\texttt{sequential}. Otherwise it is \texttt{irregular}. With a full five-read
window, this means all four transitions must match; a single non-contiguous
read is enough to flip the label. \texttt{unknown} is returned if fewer than two reads have been recorded,
since no transition can be computed yet.

\section{Confidence Scorer}
\label{sec:design-confidence}

The confidence scorer produces a value in $[0, 1]$ representing the
fraction of sequential transitions in the current window:
\[
\text{confidence} = \frac{\text{sequential transitions}}{\text{total transitions}}
\]
A score of 1.00 indicates a perfectly sequential stream; 0.00 indicates
a fully random one; values in between reflect mixed patterns. Unlike
hardware prefetchers that rely on performance counters, this score requires
no training data and is computed directly from the current window on every
read.

The adaptive mode uses a threshold of 0.7 to decide whether to prefetch.
Because prefetching also requires the phase to be \texttt{sequential}, and
the phase detector itself requires all four transitions to match, the
effective gate is all four transitions being sequential. A fully sequential
window produces a confidence of 1.0, which always clears the 0.7 threshold.
Under current parameters the confidence threshold does not add filtering
beyond what the phase check already enforces. The confidence value is logged
on every read and included in the evaluation results, so it can be compared
directly with observed overhead.

\section{Prefetching Strategies}
\label{sec:design-strategies}

JazzyFS implements three prefetching strategies, selectable at mount time
via the \texttt{JAZZYFS\_MODE} environment variable:

\begin{itemize}
    \item \textbf{Null Prefetcher} (\texttt{none}): never issues any
    read-ahead. Isolates FUSE interposition cost from prefetching cost.

    \item \textbf{One Block Lookahead} (\texttt{baseline}): unconditionally
    issues a speculative read of 4~KB immediately following every
    \texttt{read()} call, regardless of access pattern.

    \item \textbf{Feedback-Directed Sequential Prefetcher}
    (\texttt{adaptive}): issues a prefetch if and only if the current phase
    is \texttt{sequential} and confidence $\geq 0.7$. Because
    \texttt{sequential} requires all four transitions to match, all four
    must be sequential before prefetching is enabled.
\end{itemize}

The prefetch depth, set via \texttt{JAZZYFS\_PREFETCH\_DEPTH} (default 1),
controls how many 4~KB blocks are read ahead per prefetch. The three
strategies span the full range from no prefetching (\texttt{none}) to
always prefetching (\texttt{baseline}) to prefetching only when the pattern
is sequential (\texttt{adaptive}), allowing direct comparison across all
conditions.

\section{Sonification Engine}
\label{sec:design-sonification}

The sonification engine maps three internal state values to three
independent musical parameters:

\begin{itemize}
    \item \textbf{Mode} $\rightarrow$ \textbf{scale}: \texttt{none} uses
    Major; \texttt{baseline} uses Natural Minor; \texttt{adaptive} uses
    Harmonic Minor.
    \item \textbf{Phase} $\rightarrow$ \textbf{tempo}: sequential phases
    play at 0.15~s per note; irregular phases play at 1.2~s per note.
    \item \textbf{Confidence} $\rightarrow$ \textbf{melodic direction}:
    confidence $\geq 0.5$ produces an ascending melody; confidence $< 0.5$
    produces a descending melody.
\end{itemize}

A background thread detects when a workload finishes by waiting for 1
second of read inactivity, then waits an additional 5 seconds before
beginning audio playback. The phase history
is divided into segments, each played at its corresponding tempo with chord
progressions for support. Because each parameter maps to a separate musical
dimension, the design aims to make mode, phase, and confidence
perceptually distinguishable without them interfering with each other.

\section{Read Callback Logic}
\label{sec:design-read}

The core prefetching and phase detection logic executes within the \texttt{read()} callback; sonification playback runs in a separate background thread.
Listing~\ref{lst:read} shows the implementation of this callback.

\begin{lstlisting}[language=Python, caption={JazzyFS \texttt{read()} callback.}, label={lst:read}]
def read(self, path, size, offset, fh):
    # 1. Perform the actual read; confirm position with a second lseek.
    os.lseek(fh, offset, os.SEEK_SET)
    actual_offset = os.lseek(fh, 0, os.SEEK_CUR)
    data = os.read(fh, size)
    self._log_read(path, actual_offset, len(data))  # write to access.csv

    # 2. Under the state lock: reset histories on idle, append event,
    #    detect phase for sonification, update last_read_time.
    with self._state_lock:
        if time.time() - self.last_read_time > 2.0:
            self.trace.clear()
            self.phase_history.clear()
            self.confidence_history.clear()
            self.prefetch_history.clear()
        self.trace.append({"t": time.time(), "path": path,
                           "offset": actual_offset, "size": len(data)})
        sonification_phase = self._detect_phase()
        if sonification_phase != "unknown":
            self.phase_history.append(sonification_phase)
        self.last_read_time = time.time()

    # 3. Compute confidence outside lock (reads trace, no modification).
    phase      = sonification_phase
    confidence = self._prediction_confidence()

    # 4. Decide whether to prefetch.
    prefetch = False
    prefetch_offset = None
    if self.prefetch_enabled and len(data) > 0:
        if self.mode == MODE_BASELINE:
            prefetch = True
        elif self.mode == MODE_ADAPTIVE:
            prefetch = (phase == "sequential" and
                        confidence >= CONFIDENCE_THRESHOLD)
        if prefetch:
            prefetch_offset = actual_offset + len(data)
            self._prefetch_next(self._full(path),
                                prefetch_offset, self.prefetch_size)

    # 5. Update sonification histories.
    with self._state_lock:
        self.confidence_history.append(confidence)
        self.prefetch_history.append(prefetch)

    # 6. Log decision to decisions.csv.
    self._log_decision(path, actual_offset, len(data),
                       phase, confidence, prefetch, prefetch_offset)
    return data
\end{lstlisting}

The callback performs the actual read first, then runs the phase analysis
and prefetch decision. The prefetch helper reads ahead by
\texttt{prefetch\_size} bytes, which the kernel retains in the page cache
so future accesses to that data are served from memory.

\section{Implementation}
\label{sec:design-impl}

JazzyFS is a single Python file (\texttt{jazzyfs\_min.py}, approximately
880 lines) using \texttt{fusepy} for FUSE bindings and \texttt{sox} for
audio. Mode and workload name are passed in via environment variables.
Each read is logged to \texttt{access.csv} and each prefetch decision to
\texttt{decisions.csv}. Both files flush after every row so logs are not
lost if the process is interrupted.

A lock (\texttt{\_state\_lock}) protects shared state between the FUSE
callbacks and the background audio thread. Because \texttt{nothreads=True}
serializes all FUSE callbacks, no extra locking is needed inside the read
path itself.

\subsection{Experiment Harness}

The experiment harness is a separate script that automates all experimental
runs. For each combination of workload, mode, and repetition index, it:
\begin{enumerate}
    \item Sets the environment variables for that run.
    \item Mounts JazzyFS and waits for the mountpoint to become accessible.
    \item Records the start time, executes the workload script, and records
    the end time.
    \item Unmounts JazzyFS and waits for all log writes to flush.
    \item Appends the timing result to a per-run CSV file.
\end{enumerate}

After all runs complete, a separate analysis script reads the per-run CSVs,
computes summary statistics, and generates the figures and tables used in
this thesis.

\subsection{Parameter Summary}

Table~\ref{tab:params} summarizes the key design parameters of JazzyFS and
their default values.

\begin{table}[H]
\centering
\caption{JazzyFS design parameters and default values. Parameters shown in uppercase are configurable via environment variables; the remainder are hardcoded constants.}
\label{tab:params}
\begin{tabular}{lll}
\toprule
Parameter & Default & Description \\
\midrule
\texttt{JAZZYFS\_MODE}            & \texttt{adaptive} & Prefetching algorithm \\
\texttt{JAZZYFS\_PREFETCH\_DEPTH} & 1                 & Blocks per prefetch \\
\texttt{PHASE\_WINDOW}            & 5                 & Reads in classifier window \\
\texttt{CONFIDENCE\_THRESHOLD}    & 0.7               & Minimum confidence to prefetch \\
\texttt{prefetch\_size}           & 4096 bytes        & Prefetch unit size \\
\bottomrule
\end{tabular}
\end{table}

\end{document}
